% =====================================================================
\begin{thebibliography}{99}
\addcontentsline{toc}{chapter}{Bibliografía}
\markboth{Bibliografía}{}
\small

\bibitem{openstax}
J.~Abramson.
\newblock \emph{College Algebra}, 2.ª ed.
\newblock OpenStax, Rice University, 2021.
\newblock \url{https://openstax.org/details/books/college-algebra-2e}

\bibitem{blanchard}
P.~Blanchard, D.~J. Higham y N.~J. Higham.
\newblock Accurate computation of the log-sum-exp and softmax functions.
\newblock Preprint, 2019.
\newblock \url{https://arxiv.org/abs/1909.03469}

\bibitem{csapp}
R.~E. Bryant y D.~R. O'Hallaron.
\newblock \emph{Computer Systems: A Programmer's Perspective}, 3.ª ed.
\newblock Pearson, 2016.
\newblock Cap.~2, «Representing and Manipulating Information». Texto de los cursos de sistemas de Carnegie Mellon.

\bibitem{ct}
T.~M. Cover y J.~A. Thomas.
\newblock \emph{Elements of Information Theory}, 2.ª ed.
\newblock Wiley, 2006.

\bibitem{goldberg}
D.~Goldberg.
\newblock What every computer scientist should know about floating-point arithmetic.
\newblock \emph{ACM Computing Surveys}, 23(1):5--48, 1991.

\bibitem{gkp}
R.~L. Graham, D.~E. Knuth y O.~Patashnik.
\newblock \emph{Concrete Mathematics: A Foundation for Computer Science}, 2.ª ed.
\newblock Addison-Wesley, 1994.
\newblock Cap.~2, «Sums».

\bibitem{hammack}
R.~Hammack.
\newblock \emph{Book of Proof}, 3.ª ed.
\newblock Richmond (Virginia), 2018.
\newblock \url{https://richardhammack.github.io/BookOfProof/}

\bibitem{ieee754}
IEEE.
\newblock \emph{IEEE Standard for Floating-Point Arithmetic} (IEEE Std 754-2019).
\newblock IEEE, 2019.

\bibitem{mcs}
E.~Lehman, F.~T. Leighton y A.~R. Meyer.
\newblock \emph{Mathematics for Computer Science}, revisión del 6 de junio de 2018.
\newblock Massachusetts Institute of Technology. Texto del curso 6.042J.
\newblock \url{https://courses.csail.mit.edu/6.042/spring18/mcs.pdf}

\bibitem{mit6042}
A.~R. Meyer y A.~Chlipala.
\newblock 6.042J: Mathematics for Computer Science.
\newblock Massachusetts Institute of Technology, MIT OpenCourseWare, primavera de 2015.
\newblock \url{https://ocw.mit.edu/courses/6-042j-mathematics-for-computer-science-spring-2015/}

\bibitem{spivak}
M.~Spivak.
\newblock \emph{Calculus}, 4.ª ed.
\newblock Publish or Perish, 2008.
\newblock Cap.~1, «Basic Properties of Numbers», y cap.~2, «Numbers of Various Sorts».

\end{thebibliography}
